%% ma: L10-B15-0046
%% lop: 10
%% chuong: 6
%% bai: 15
%% dang: 10.15.4
%% loai: TL
%% muc: TH
%% dap_an: "-"
%% nguon: "Ảnh giáo viên gửi 21/11/2026: Chuyên đề VI – Toán 10 – Chương VI (Hàm số – đồ thị và ứng dụng), tr. 199–203, Dạng 4 Câu 1"
%% kien_thuc: "Lớp 10 Bài 15 (hàm số đồng biến, nghịch biến)"
%% trang_thai: chinh-thuc
%% ghi_chu: "Dạng do giáo viên cung cấp (21/11/2026)."
\begin{ex}
Xét tính đồng biến và nghịch biến của hàm số $f(x)=x^2-7$ trên khoảng $(-\infty;0)$ và trên khoảng $(0;+\infty)$.
\loigiai{Với $x_1\ne x_2$: $\dfrac{f(x_1)-f(x_2)}{x_1-x_2}=\dfrac{x_1^2-x_2^2}{x_1-x_2}=x_1+x_2$.

Trên $(-\infty;0)$: $x_1+x_2<0$ nên hàm số nghịch biến. Trên $(0;+\infty)$: $x_1+x_2>0$ nên hàm số đồng biến.}
\end{ex}
